\documentclass[11pt]{article}

\usepackage[margin=1in]{geometry}
\usepackage{amsmath,amssymb,amsthm}
\usepackage{mathtools}
\usepackage{xcolor}
\usepackage{enumitem}
\usepackage{tcolorbox}
\usepackage{hyperref}
\hypersetup{colorlinks=true, linkcolor=black, urlcolor=blue}

% ---- status tags ----
\newcommand{\est}{\textcolor{teal!70!black}{\textsf{\small[EST]}}}
\newcommand{\fields}{\textcolor{violet}{\textsf{\small[FIELDS]}}}
\newcommand{\asm}{\textcolor{blue!70!black}{\textsf{\small[ASM]}}}
\newcommand{\conj}{\textcolor{orange!80!black}{\textsf{\small[CONJ]}}}
\newcommand{\post}{\textcolor{purple!70!black}{\textsf{\small[POST]}}}
\newcommand{\gap}{\textcolor{red!75!black}{\textsf{\small[GAP]}}}
\pdfstringdefDisableCommands{%
  \def\est{[EST]}\def\fields{[FIELDS]}\def\asm{[ASM]}%
  \def\conj{[CONJ]}\def\post{[POST]}\def\gap{[GAP]}%
  \def\quad{ }\def\normalsize{}%
}

% ---- theorem environments ----
\theoremstyle{definition}
\newtheorem{definition}{Definition}[section]
\newtheorem{proposition}{Proposition}[section]
\newtheorem{claim}{Claim}[section]
\newtheorem{corollary}{Corollary}[section]
\newtheorem{openproblem}{Open Problem}[section]

\theoremstyle{remark}
\newtheorem*{remark}{Remark}
\newtheorem*{invariant}{Invariant I (necessity of partition)}

\newtcolorbox{principlebox}{colback=gray!6, colframe=gray!40, boxrule=0.5pt, arc=2pt, left=6pt, right=6pt, top=5pt, bottom=5pt}

\title{\textbf{A Formal Skeleton for the Cut}\\[4pt]
\large An operational core, one quantum realization, and explicit limits of inference}
\author{}
\date{Working spec, v0.5 --- 2026-09-23}

\begin{document}
\maketitle

\begin{principlebox}
\noindent\textbf{Status.} This is a skeleton, not a finished theory. Every claim is tagged:
\begin{itemize}[leftmargin=1.4em, itemsep=1pt, topsep=3pt]
  \item \est{} --- established result in the literature, cited, not original here.
  \item \fields{} --- due to Chris Fields (and collaborators); this framework assembles it, does not originate it.
  \item \asm{} --- assembly: a combination or restatement of established pieces, possibly original in \emph{arrangement}, using no new mathematics.
  \item \conj{} --- conjecture: plausible, not proved here, flagged as open.
  \item \post{} --- postulate: a philosophical commitment of the framework's ontology, not a physics result; adopted, not derived.
  \item \gap{} --- a deliberately stated limitation, missing bridge, or scope boundary. It is a no-go result only where a separate proof is supplied.
\end{itemize}
\vspace{2pt}
\noindent\textbf{Design principle.} Keep the proposed operational core, particular physical realizations, and downstream phenomena or interpretations distinct. Formalize the limits of each claim rather than turning an assumption of the realization into observer-accessible evidence. The scope of the partition-selection argument is stated in \S\ref{sec:select}.

\vspace{2pt}
\noindent\textbf{Scope of this realization.} Given a chosen cut and declared state, symmetry, interaction, and apparatus, the model defines subsystem entropy, cross-cut flux, and observer-accessible records. The present input tuple does not select that cut or derive a macroscopic arrow. These model outputs do not uniquely determine ontology or rule out a richer physical cut-selection principle.
\end{principlebox}

\section*{Preface: operational core and scope}
\addcontentsline{toc}{section}{Preface: operational core and scope}

The canonical statement of the new module is \href{../docs/operational-core.md}{\texttt{docs/operational-core.md}}; its public, static presentation is \href{../web/operational-core.html}{\texttt{web/operational-core.html}}. It is a \textbf{proposed foundation for this research programme}, not a new definition of time, not a replacement for the ontology of The Cut, and not a claim that the components are the unique minimal or fully logically independent axioms.

The framework now has three layers. The \textbf{operational core} specifies procedures, records, comparison rules, and the scope of inference, but supplies neither dynamics nor an ontology of the whole. \textbf{Physical realizations} add structures that could implement those capacities and must declare their additional assumptions. \textbf{Phenomena and interpretations} study clocks, arrows, records, shareability, and experience within a realization; they do not follow from the operational core alone. The layers are related by adding assumptions and testing consequences, not by automatic deduction. A realization or interpretation may fail without the entire operational core failing.

The operational basis comprises distinguishability (OC-A1), recordability and comparability (OC-A2), composability of some preparation--interaction--readout procedures (OC-A3), and restricted regularity within a stated domain (OC-A4). OC-A3 openly incurs a foundational debt: composition and retained traces already assume a primitive local order. The core therefore does \emph{not} derive all of time from a wholly order-free basis. This dated clarification (2026-09-23) separates causal accessibility and ordering from calibrated duration, history and record asymmetry, and temporal judgement; it does not revise historical experiment records. Its inference constraints are predictive consistency for complete accessible transcripts (OC-R1), invariance under mere changes of representation for the same physical procedure (OC-R2), and scope-limited identifiability relative to the available procedure class $\mathcal Q_O$ (OC-R3).

This document supplies one \textbf{physical realization}: a chosen tensor factorization, global pure state, unitary dynamics, and a declared Hamiltonian split, together with the apparatus assumptions used in particular models. These are model assumptions, not consequences of OC-A1--A4 and not direct observations of the whole universe. A model cut is likewise distinct from operational access: the former partitions the modeled system and interactions; the latter is the observer's actual preparations, controls, readout channels, records, resolutions, and resources, which generate $\mathcal Q_O$. The mapping between them is neither assumed to be one-to-one nor specified by the core.

An observer here may be any physical record-capable system; consciousness is not assumed. Nor does it receive a privileged model-builder view. Any endogenous test must declare its clock, memory, phase reference, preparation source, and signal channels. Simulator steps, the full state matrix, and debug variables are not automatically observer records.

\section{Primitives and the one invariant}\label{sec:prim}

For the quantum realization studied in this spec, let $U$ be a closed quantum system with Hilbert space $\mathcal{H}_U$, assumed global pure state $\lvert\Psi\rangle$, and total Hamiltonian $\hat{H}_U$, assumed to evolve unitarily. These are \textbf{model assumptions}. Standard quantum mechanics establishes the consequences conditional on them; this spec does not claim observational evidence that the universe as a whole actually satisfies them.

A \textbf{cut} is a choice of tensor factorization
\[
  \mathcal{H}_U \;\cong\; \mathcal{H}_A \otimes \mathcal{H}_B
\]
together with a designation of $A$ as ``subsystem'' and $B$ as ``complement.'' \fields{} (Fields calls the interface between them a \emph{holographic screen}; the factorization is the formal content.)

\begin{definition}[trivial cut]
The trivial cut is $A = U,\ B = \mathbb{C}$ (or $B = U,\ A = \mathbb{C}$): no genuine bipartition.
\end{definition}

\begin{invariant}[\asm]
The \emph{cut-relative} indicators used here (reduced entropy, cross-cut flux, and cross-cut registration) degenerate on the trivial cut: no complement exists to trace out, no cross-cut boundary exists for flux, and no complement channel exists for registration. A genuine bipartition makes these indicators definable but does not guarantee nonzero values; those depend on state, dynamics, and access. This statement concerns cut-relative diagnostics, not whether causal accessibility can exist within a model without a chosen bipartition.
\end{invariant}

\noindent The necessity of partition is primitive but silent: it is invariant (it holds for every genuine cut) and yet it does not select among them (see \S\ref{sec:select}).

\section{Causal accessibility, clocks, and the arrow \quad{\normalsize\asm\ / \conj}}\label{sec:time}

This section adds a model-relative ordering layer to the operational core. It uses the primitive local sequencing already admitted by OC-A2--A3; it does not derive directed order from an order-free substrate. The chosen cut, state, dynamics, admissible interventions, and readouts are additional data of the quantum realization. Neither a correlation nor a nonzero Hamiltonian matrix element alone proves directed influence.

\subsection*{Causal accessibility and scoped order}
For declared events or operational regions $e,f$, write $e\leadsto_M f$ when a controlled change associated with $e$ can alter the distribution of an accessible record at $f$, with other declared conditions fixed. Schematically,
\[
 e\leadsto_M f\quad\Longleftrightarrow\quad
 \exists x,x',Q,r_f:\ P_M(r_f\mid\operatorname{do}(x),Q)\ne P_M(r_f\mid\operatorname{do}(x'),Q).
\]
The intervention/preparation alternatives $x,x'$ belong to $e$; $Q$ fixes the rest of the procedure; $r_f$ belongs to $f$. This is a \emph{model-level influence relation}, not a claim that every observer can test the edge. Common-source correlation does not establish $e\leadsto_M f$ or its reverse.

The observer-identifiable relation $\leadsto_{\mathcal Q_O}$ is reconstructed only from interventions, detectors, clocks, phase references, records, resolution, and trial budget in the access class $\mathcal Q_O$. In general $\leadsto_{\mathcal Q_O}\ne\leadsto_M$: an edge can exist in the model without being identifiable to this observer, and finite, inadequately controlled data can also support a false edge. The model cut is not the access class. The reconstruction problem is how much of $C_M=(E,\leadsto_M)$ can be recovered as $\widehat C_O$ from accessible transcripts. The 2026-09-22 targeted follow-up illustrates this scope distinction in the same modeled sample and dynamics: the original battery has an approximately quadratic weak-coupling orientation response, while an added pair-sensitive local reference and joint readout give an approximately linear response; the zero-coupling null remains. These are additional apparatus resources, so detectable leading order changes with the access class. The result does not imply that the sample's physical causal structure changes with measurement choice.

Define $e\prec_M f$ by a directed path of one or more $\leadsto_M$ edges. If the relevant domain is acyclic, this transitive reachability is a strict partial order (and its reflexive closure is a partial order). If loops are permitted it remains directed reachability, without automatic antisymmetry or irreflexivity. Acyclicity requires an assumption or a demonstration; unrelated events need not be placed in a universal line. This order carries no duration by itself.

\subsection*{Clock relations and duration}
A clock is a declared physical reference and read/count procedure with a calibration domain and matching rule. Clock relations can assign duration along a causal chain or worldline only with that extra structure. Keep the model parameter $t$, causal reachability $\prec_M$, physical clock readings, relativistic proper time in its applicable domain, and any entropic variable distinct. No universal seconds or proper-time law follows from reachability alone.

\subsection*{Arrow, records, and ``time as forgetting''}
Given a cut $A\vert B$ and global state $\rho=\lvert\Psi\rangle\langle\Psi\rvert$, the reduced state and its von Neumann entropy are
\[
 \rho_A=\operatorname{Tr}_B(\rho),\qquad S_A=-\operatorname{Tr}(\rho_A\log\rho_A).\qquad\est
\]
Boundary/preparation conditions, environment, coarse-graining, record formation and persistence, loss of accessible information, and recoverability are relevant to a historical arrow. Entropy increase cannot be both the definition of that arrow and its explanation. The conditional phrase ``time as forgetting'' now concerns asymmetric retention or loss of causal records \emph{along an already ordered history}; it is not the origin of causal order.

Write $\tau_A^{\mathrm{ent}}$ only as a candidate entropic-orientation coordinate on a declared segment where $S_A$ is sufficiently monotone. It is model- and cut-dependent, can fail globally, and supplies neither a universal duration nor proper time. A causal chain $e_1\prec_M e_2\prec_M e_3$ can persist when $S_A(e_1)=S_A(e_3)$ or entropy temporarily falls. A recurrence invalidates a universal entropic clock, not the causal chain.

\begin{proposition}[global purity, local mixedness]{\normalfont\ \est}\label{prop:purity}
If $\lvert\Psi\rangle$ is pure, then $S_U=0$, while $S_A\ge0$ with equality iff $A\vert B$ is unentangled. An entangled part has $S_A>0$; this does not define causal order, duration, or a macroscopic arrow. A varying $S_A$ can label an entropic orientation only where monotonicity is established for the specified state, cut, and interval. Global purity does not imply stationarity or the absence of internal temporal relations.
\end{proposition}

\begin{remark}[direction and records]{\normalfont\ \gap}
Unitarity supplies reversible evolution, while time-reversal symmetry is a separate dynamical condition. A state and dynamics determine the sign of a specified entropy change, but the entropy functional alone supplies no universal macroscopic arrow. A physical record need not always raise the recorder's reduced entropy. Preparation, correlations, environment, erasure, and resources must be modelled; any claimed arrow also needs a history-level test rather than a definition that guarantees its own result.
\end{remark}

\noindent\textbf{Four distinct tasks.} Ordering asks which interventions can affect which records; duration asks how calibrated clocks compare along those relations; arrow asks why histories and records are statistically asymmetric; temporal judgement asks how a system integrates signals and memories to infer before/after, duration, and ``now.'' The last task requires a model of integration, prediction, belief update, and report, but is not required for physical order to exist. None is an automatic deduction from the others.

\section{Energy as broken symmetry \quad{\normalsize\est}\label{sec:energy}}

By Noether's theorem, energy is the conserved charge of time-translation symmetry. \est{}

Given a cut \textbf{and a chosen local Hamiltonian decomposition}
\[
  \hat{H}_U \;=\; \hat{H}_A\otimes I_B + I_A\otimes\hat{H}_B + \hat{H}_{AB}, \qquad \fields
\]
the interaction term $\hat{H}_{AB}$ defines boundary exchange. \textbf{The cut makes such a decomposition possible; it does not uniquely select it} --- the split into $\hat{H}_A$, $\hat{H}_B$, $\hat{H}_{AB}$ is a further choice, on top of the factorization, exactly as Claim~\ref{cl:noselect} says the factorization itself is a choice on top of $\mathcal{H}_U$.

\begin{proposition}[local energy flux]{\normalfont\ \est}\label{prop:eflux}
For a fixed split $\hat{H}_U = \hat{H}_A + \hat{H}_B + \hat{H}_{AB}$ with $\hat{H}_A$ time-independent and $[\hat{H}_A,\hat{H}_B]=0$, the instantaneous energy flux into $A$ is
\[
  \dot{E}_A \;=\; \frac{i}{\hbar}\big\langle [\hat{H}_{AB}, \hat{H}_A]\big\rangle.
\]
Hence $\hat{H}_{AB}=0$ is \emph{sufficient} for no exchange, but not necessary: local energy is generically non-conserved when the interaction fails to commute with the local Hamiltonian, whereas a coupling that commutes with $\hat{H}_A$ (or a special state) can give zero flux even with $\hat{H}_{AB}\neq 0$. Energy is a cut-relative quantity that generically crosses the cut, without ``$\hat{H}_{AB}\neq 0$'' alone being the criterion.
\end{proposition}

\begin{proposition}[whole has no energy-flow]{\normalfont\ \asm}
On the trivial cut there is no $\hat{H}_{AB}$, hence no energy exchange; the ``energy'' that flows is strictly cut-relative. The global $\hat{H}_U$ conserves; only a part leaks.
\end{proposition}

\begin{remark}[\asm]
In Fields' construction, the bits exchanged across the screen are proportional to the transferred energy represented by the $\hat{H}_{AB}$ eigenvalue. The relation holds at the declared boundary; it is not an identity between energy and information in general.
\end{remark}

\section{Information as registered difference \quad{\normalsize\est}\label{sec:info}}

Two layers, kept separate.

\paragraph{4a. Shannon layer. \est}
For a source with distribution $\{p_i\}$, information content is $-\log p_i$, and average surprise is $H = -\sum_i p_i \log p_i$. This is formally parallel to \S2's $S_A$ (von Neumann entropy equals Shannon entropy on the eigenvalue distribution of $\rho_A$). The shared form does \textbf{not} identify uncertainty, information transfer, physical duration, and record order. Connecting those roles requires a model of preparation, channel, readout, and clock calibration.

\paragraph{4b. Observer-registration layer.}
Not all difference is information \emph{to a system}. Given an observer channel $K_B : \Omega \to \Omega_B$, a difference between source states $A, A'$ is \textbf{registered} iff
\[
  K_B(A) \;\neq\; K_B(A'),
\]
and registered \emph{above threshold} iff $d_B\big(K_B(A), K_B(A')\big) \ge \delta_B$.

\begin{proposition}[no complement-relative registration on the trivial cut]{\normalfont\ \asm}
A single global pure state has zero von Neumann entropy, and the trivial cut offers no \emph{complement}-relative distribution --- there is no $B$ to trace out and no cross-cut channel. (The global system may still contain internal subsystems, Born distributions, clocks, and records relative to further structure.) What is absent is the complement-relative registered information and reduced-entropy diagnostic defined for this trivial cut, not causal order, information, or time as such.
\end{proposition}

\begin{remark}[\gap\ --- which channel]
$K_B$ is \emph{given}, not derived. Which channel an observer has --- its threshold $\delta_B$, its representation space $\Omega_B$ --- is an input. Note the parallel to \S\ref{sec:select}: just as the tensor factorization is chosen, the channel is chosen. Same gap, observer-side.
\end{remark}

\begin{remark}[recording is not a universal entropy-increase rule]
That a physical interaction leaves a usable trace does not imply that every recording operation increases the reduced-state entropy of the observer. The answer depends on preparation, correlations, environment, erasure, and resource accounting. Record order and entropy change may be related in a specified realization, but are not identical by definition.
\end{remark}

\section{Reach: what else the cut defines, and what it does not}\label{sec:reach}

The cut-relative quantities discussed in \S2--\S4 are not the framework's ceiling. Four routes carry further physical quantities out of the same apparatus --- and one route provably does not run, which is as much a result as the others.

\paragraph{Route 1 --- Noether charges, by the energy move. \est}
Energy (\S3) was the conserved charge of \emph{time}-translation symmetry, made cut-relative through the boundary term $\hat{H}_{AB}$. Noether's theorem is general: every continuous symmetry yields a conserved charge, and — with the same caveat as the energy flux — each splits across the cut.

\begin{proposition}[charge flux across the cut]{\normalfont\ \est}
Energy is special: the Hamiltonian both \emph{is} the time-translation charge and \emph{generates} the dynamics. A generic Noether charge $\hat{Q}$ does not generate evolution. Its flux is generated by the \emph{Hamiltonian}:
\[
  \dot{Q}_A \;=\; \frac{i}{\hbar}\big\langle [\hat{H}_U, \hat{Q}_A]\big\rangle,
\]
and, when a local current $j_Q$ exists, equals a boundary flux $\dot{Q}_A = -\int_{\partial A} j_Q\!\cdot n$. So: for global symmetries with local current decompositions, subsystem charge changes are boundary fluxes generated by the Hamiltonian dynamics; $\hat{Q}_U$ may be globally conserved while $\hat{Q}_A$ is not. This covers momentum (space-translation), angular momentum / spin (rotation), and other global charges --- \emph{without} forcing each into the same generator role as energy. Energy is the case where the charge and the generator coincide, which is why it was examined first.
\end{proposition}

\begin{remark}[gauge charges are subtler]{\normalfont\ \est\ / \gap}
Electric and other \emph{gauge} charges do \textbf{not} automatically fit the naive $\mathcal{H}_A\otimes\mathcal{H}_B$ split. Gauge constraints (Gauss's law) can obstruct simple Hilbert-space factorization across a spatial region; a gauge-invariant treatment requires local algebras, edge modes, or an extended Hilbert space (Casini--Huerta--Rosabal). Route~1 as stated applies cleanly to global symmetries with local currents; gauge symmetries need the algebraic/edge-mode version and are flagged here rather than claimed.
\end{remark}

\noindent This extends the construction from three listed quantities to the conserved charges associated with the symmetries already assumed by the model.

\paragraph{Route 2 --- thermodynamic quantities, from the reduced entropy. \est\ (cut + arrow)}
With $S(\rho_A)$ (\S2) and $\langle\hat{H}_A\rangle$ (\S3) in hand, standard relations deliver the thermodynamic quantities. Temperature is the marginal exchange rate between the information-face and the energy-face of the same cut,
\[
  \frac{1}{T} \;=\; \frac{\partial S(\rho_A)}{\partial \langle\hat H_A\rangle},
\]
and free energy, chemical potential, and pressure follow by the usual derivatives. \textbf{Caveat:} these relations presuppose equilibration and the relevant thermodynamic conditions. They therefore require both a subsystem choice and additional arrow conditions, unlike the Noether charges, which require the subsystem decomposition but not a thermodynamic arrow.

\paragraph{Route 3 --- invariants and geometry. \est\ / \conj}
Two sub-cases of different status:
\begin{itemize}[leftmargin=1.4em, itemsep=2pt, topsep=2pt]
  \item \textbf{Mass. \est} Rest mass is a Casimir invariant of the Poincar\'e group, $m^2 = E^2 - p^2$: \emph{frame}-invariant within a chosen subsystem representation. Given energy and momentum from Route~1, mass is their frame-invariant combination --- the feature of the energy--momentum content that does not vary under boosts. Note this is Lorentz-invariance, not cut-invariance: the mass/energy assigned to a subsystem still varies with binding, dressing, and boundary choice.
  \item \textbf{Distance / geometry. \conj} In holographic settings (AdS/CFT), Ryu--Takayanagi relates the entanglement entropy of a boundary region to the area of a bulk extremal surface --- a specific relation, not a general ``mutual entanglement $=$ distance'' rule. More broadly, entanglement/geometry programs (Van~Raamsdonk) suggest connectivity and distance-like structure may be recoverable from entanglement patterns, with disentangling regions pinching spacetime apart. This framework is \emph{compatible} with that program but does not derive it, and does not assert the general rule.
\end{itemize}

\paragraph{Route 4 --- what the cut does NOT define. \gap}
Not every physical quantity is a co-product of a cut, and marking the exclusions is a result, not an omission.
\begin{itemize}[leftmargin=1.4em, itemsep=2pt, topsep=2pt]
  \item \textbf{The dimensionful constants} $(\hbar, G, \text{coupling constants})$ are fixed background: identical across every partition and observer, cut-\emph{invariant} in the strongest sense, and \emph{not generated} by the cut. (The status of $c$ is subtler and is treated in \S\ref{sec:lorentz}: as the observed relativistic \emph{conversion constant} between space- and time-units it is a convention/background scale, but an \emph{emergent maximum correlation speed} of a different, pre-geometric character may arise from the interaction structure. These are two distinct roles for the one symbol and must not be conflated.)
  \item \textbf{The symmetry group and dimensionality} (which gauge group, why $3{+}1$ dimensions) are \emph{upstream} of the cut: Route~1 \emph{consumes} the symmetries as input and cannot produce them. Likewise the spacetime dimension.
\end{itemize}

\begin{proposition}[scope of the current cut construction]{\normalfont\ \asm}\label{prop:typebound}
Given the declared background structure, the current realization defines partition-relative charges, correlations, exchange terms, and related thermodynamic quantities. It does \textbf{not} derive the constants, symmetry group, dimensionality, or the partition itself. Proposed geometry-from-entanglement links require additional assumptions and are not consequences of the cut alone.
\end{proposition}

\paragraph{Matter: the conserved side of the cut. \asm}
The quantities so far include flows and changing correlations: energy may cross $\hat{H}_{AB}$, records may be created or lost, and $S_A$ may label an entropic orientation on a restricted interval. Matter is interpreted below as the \emph{persisting} side; this is a model-dependent proposal, not a definition of time.

\begin{definition}[matter]\label{def:matter}
The \textbf{matter} content of a subsystem $A$ is the conserved, boundary-holding structure of $A$ --- the invariant of $A$'s own evolution: what is preserved (protected by symmetry, commuting with the effective $\hat{H}_A$) while the correlations that decohere into $B$ are washed away. Where energy is what flows across the cut, matter is what stays.
\end{definition}

\noindent This is not a definition in terms of energy located in $\hat{H}_A$. It is an interpretive classification of structures that remain stable over the modeled interval. It does not derive matter from entropy growth.

\begin{proposition}[consequences and limits of the matter interpretation]{\normalfont\ \asm}\label{prop:matter}
Within a fixed representation and operational regime, Definition~\ref{def:matter} has the following consequences and limits:
\begin{itemize}[leftmargin=1.4em, itemsep=1pt, topsep=2pt]
  \item \textbf{Inertia / persistence} is not added; ``conserved under evolution'' is what resisting change \emph{means}.
  \item \textbf{Localization}: a conserved charge in a region is a boundary --- it is what stays on $A$'s side while correlations leak.
  \item \textbf{Trace-retention}: matter is what can register (O2, Def.~\ref{def:observer}). Strictly, a detector requires not exact conservation but a robust \emph{pointer subspace/subalgebra} whose records persist over the relevant timescale (decoherence / quantum Darwinism); ``persistent structure'' should be read in that sense, not as strict invariance.
  \item \textbf{Mass coincidence}: by Route~3a, rest mass is the Poincar\'e Casimir --- \emph{frame}-invariant within a chosen subsystem representation. (It is not straightforwardly \emph{cut}-invariant: the mass/energy assigned to a subsystem varies with binding, dressing, and boundary choice.) With that scoping, mass being the \emph{frame-invariant} combination and matter being the \emph{persistent} aspect are aligned --- a consistency check, once the representation is fixed.
\end{itemize}
\end{proposition}

\begin{remark}[three flags]{\normalfont\ \gap}
(i)~\textbf{Cut-relative.} ``Conserved under $\hat{H}_A$'' needs an evolution, hence a cut; change the cut and different structure is conserved. There is \emph{no primitive matter} --- only cut-relative persistence. Matter is not more fundamental than the cut.
(ii)~\textbf{Arrow dependence of this interpretation.} Reading persistence as an asymmetric historical notion requires the boundary, environmental, and coarse-graining assumptions used for the arrow. Conservation and stable correlations can still be defined in reversible dynamics; the stronger claim that an arrow is the condition for anything to exist is not inferred.
(iii)~\textbf{Category, not content.} This defines matter-as-a-category (conserved boundary-holding structure), not the specific spectrum (why these particles). The particle content, like the symmetry group, is Route~4 background: consumed, not derived.
\end{remark}

\begin{remark}[flux and persistence]{\normalfont\ \asm}
As an interpretive heuristic, the spec distinguishes boundary fluxes and changing correlations from persistent structures. This classification is not a definition of matter, energy, or time. A macroscopic arrow still requires the additional conditions stated in \S\ref{sec:time}.
\end{remark}

\section{Particles and forces \quad{\normalsize\asm\ / \conj}\label{sec:pf}}

\paragraph{Particles as persistent charge sectors. \asm}
Definition~\ref{def:matter} treats matter as conserved boundary structure. The following interpretation relates particle types to persistent charge sectors.

\begin{definition}[particle]\label{def:particle}
A \textbf{particle} is a minimal persistent charge-sector \emph{relative to the theory's representation content and operational regime}: an indivisible package of frame-invariant charges (mass, spin, internal charges) that is not decomposable, within that regime, into independently-persisting sub-sectors carrying the same conserved package.
\end{definition}

\noindent Wigner's classification identifies a relativistic particle type with an irreducible unitary representation of the Poincar\'e group, labelled by its Casimirs. The interpretation above is consistent with that classification in a stated regime, but it does not cover every use of ``particle'' or derive the particle spectrum. Confinement, infraparticles, resonances, quasiparticles, and partons require their own regimes and definitions.

\paragraph{Interactions across a cut. \asm/\est}
The term $\hat{H}_{AB}$ models coupling between the two sides and may transfer locally defined charges (Prop.~\ref{prop:eflux}). This decomposition describes interaction in the realization; it is not a complete definition of force.

\paragraph{Candidate relation to gauge structure. \conj}
The spec asks whether independently chosen internal conventions across regions require a connection that relates them. Gauge theory already supplies such connections and their associated fields. The proposed link to cut-relative symmetry is interpretive: it does not derive gauge fields, gauge groups, or gauge bosons from a partition alone.

\paragraph{Status of the comparison.}
Two established settings constrain this interpretation.

\begin{itemize}[leftmargin=1.4em, itemsep=3pt, topsep=3pt]
  \item \textbf{Aharonov--Bohm is loop-relative and gauge-invariant. \est-supported} The interference depends on enclosed flux even when the particle travels through a field-free region. This supports a holonomy description, but it does not determine the ontology of the potential or explain the origin of the connection.
  \item \textbf{Gauge-theory factorization depends on algebra and boundary choices. \conj} Casini, Huerta, and Rosabal show that a gauge-invariant Hilbert space need not tensor-factorize across a region. Different algebra and edge-mode prescriptions can give different regional entanglement entropies. This motivates explicit cut and algebra choices, but does not prove the proposed gauge interpretation or select a unique cut.
\end{itemize}

\begin{remark}[limits of the gauge comparison]{\normalfont\ \conj/\gap}
(i)~Regional entanglement entropy may depend on the boundary and algebra prescription, and may diverge. Clearly defined relational quantities such as mutual information or relative entropy are therefore preferable where applicable. This ambiguity does not prove a universal no-selection claim.
(ii)~Relational quantum-reference-frame work is prior research. The spec offers an informal comparison and does not originate or derive that programme.
(iii)~The comparison does not determine the gauge group or explain why a symmetry is gauged. Both remain additional physical input.
\end{remark}

\section{Worked example: black-hole information and the island cut \quad{\normalsize\asm, consistency check}\label{sec:bh}}

This section is a \textbf{scope comparison, not a contribution to black-hole physics}. The framework does \textbf{not} derive the island rule, Page curve, or quantum extremal surface (QES) formula. Those results require the gravitational $\mathrm{Area}/4G$ term, holography, and semiclassical path-integral methods.\footnote{Standard references for the machinery cited here: Hawking (radiation, information loss); Page (the Page curve); Ryu--Takayanagi (holographic entanglement entropy, hep-th/0603001); Engelhardt--Wall / QES; Penington (entanglement-wedge reconstruction, 1905.08255); Almheiri--Engelhardt--Marolf--Maxfield and Almheiri--Mahajan--Maldacena--Zhao (island rule, 1905.08762 and follow-ups). The framework uses these results; it does not reproduce or derive them.}

\paragraph{Scope issue used in the comparison.}
A black hole forms from a pure state and evaporates. Hawking's calculation makes the radiation's entanglement entropy rise monotonically, so the final radiation looks highly mixed --- apparently destroying information, violating unitarity. In the framework's vocabulary this warns against \textbf{treating a cut-relative local mixedness as a cut-independent global fact}. The comparison with the problem of time is an analogy about scope, not a derivation or an assertion that global purity eliminates all temporal relations.

\paragraph{Island calculation.}
The island / Page-curve resolution: the radiation's fine-grained entropy rises, peaks at the Page time, and falls back --- the Page curve --- because after the Page time a region \emph{inside} the horizon (the island) enters the \emph{entanglement wedge} of the radiation. The radiation entropy is computed by the dominant generalized-entropy saddle, extremized/minimized over candidate islands:
\[
  S(R) = \min_I\;\mathrm{ext}_I\Big[\tfrac{\mathrm{Area}(\partial I)}{4G_N} + S_{\text{bulk}}(R\cup I)\Big].
\]
The Page-time transition is a \textbf{change in which saddle dominates}. The comparison uses three features:
\begin{itemize}[leftmargin=1.4em, itemsep=2pt, topsep=2pt]
  \item \textbf{Global purity and local mixedness (Prop.~\ref{prop:purity}).} Radiation entropy is a subsystem quantity. The dominant QES/island saddle gives the fine-grained entropy used in the Page curve, consistent with global unitarity in the cited models.
  \item \textbf{Cut-relativity with a chosen cut (Claim~\ref{cl:noselect}).} The physics does not merely permit cut-relative entropy; it \emph{computes the answer by extremizing over candidate islands / QES saddles}. Which saddle dominates changes over time.
  \item \textbf{Declared entropy prescription.} As in the gauge case (\S\ref{sec:pf}), the subsystem and entropy prescription must be stated before interpreting the result.
\end{itemize}

\paragraph{Limit of the comparison. \gap}
The island is a \emph{geometrically disconnected} interior region that nonetheless belongs to the radiation's wedge. The framework can call this ``a non-obvious cut,'' but it has no machinery for \emph{why} the correct cut includes a spatially separated interior --- that comes from the $\mathrm{Area}/4G$ gravitational term, which the framework does not generate (Route~4). It cannot derive the QES formula, the value of the Bekenstein--Hawking entropy, or the interior-reconstruction / firewall structure; those are holography and the gravitational path integral, consumed not explained.

\begin{remark}[scope of the analogy]{\normalfont\ \asm}
The problem-of-time discussion and the island calculation both require care about which subsystem, observable, and procedure define a quantity. This is a limited structural analogy. It is not a common derivation, a proof that the two problems have the same cause, or a prediction of the island result. Questions about interior reconstruction, state dependence, microscopic interpretation, and firewall-related issues remain outside this framework.
\end{remark}


\section{Causal structure, clocks, and spacetime \quad{\normalsize\conj\ / \gap}}\label{sec:spacetime}

The architectural guide is $\text{causal order}\to\text{causal/conformal structure}$ and then $\text{causal structure}+\text{clock, scale, or volume information}\to\text{metric structure}$. This is not a derivation of Lorentzian geometry from The Cut. In appropriate Lorentzian settings the null/causal relations constrain conformal structure, while metric scale needs additional information and assumptions. An arbitrary influence graph does not automatically become a spacetime or supply proper time.

\paragraph{Representations and physical content.}
A coordinate $t$ is a model parameter, not an embedded observer's clock. Changing a chart or allowing negative coordinate values does not reverse a physical history. Causal reachability can be represented in different coordinates without changing which events can influence which others. Physical clock comparisons add duration information along selected chains; they require calibration and a model relating clocks.

\paragraph{The ``3'' and the ``1'' are not derived here.}
Spatial labels and dimension are inherited from the realization's assumed symmetry/locality structure, not selected by the cut. Causal reachability orders some events without a universal linear coordinate. Their different modelling inputs do not prove that space and time have fundamentally different ontological types or establish a Lorentzian signature.

\begin{proposition}[distinct inputs do not establish ontological types]{\normalfont\ \asm/\gap}\label{prop:diftype}
The present realization declares spatial and causal/clock structures separately. Their distinct roles constrain a proposed model, but do not imply a fundamental type distinction between physical space and time. Neither $S_A$ nor $\tau_A^{\mathrm{ent}}$ supplies the missing causal or metric structure.
\end{proposition}

\paragraph{Penrose diagrams as a scoped example.}
A Penrose diagram is a conformal compactification of a specified spacetime model, not a literal picture. Its finite placement of infinity, drawing distances, coordinate shape, and most visual angles are representational; the null-direction convention and global relationships preserve causal reachability, horizons, and causal boundaries. The full diagram is a \emph{model-level causal completion}, not an image directly seen by an embedded observer. Such an observer has local signals, interventions, clocks, records, and other declared instruments and may reconstruct only a partial $\widehat C_O$. See, for example, \href{https://www.maths.ed.ac.uk/~jsbiersk/assets/LectureNotesGRII.pdf}{Sbierski's General Relativity II notes, \S3.2} for the conformal construction.

\paragraph{Clock and metric bridge. \conj/\gap}
A physical realization would have to establish an appropriate local causal structure, calibrated clocks, scale or volume information, and Lorentzian transformation laws before claiming metric duration or $3{+}1$ kinematics. The model's entropy may orient particular recorded histories after those histories are specified; it is not the index from which spacetime is built.

\begin{proposition}[conditional local kinematics]{\normalfont\ \asm/\conj}\label{prop:threeone}
If a realization supplies the requisite causal and Lorentzian structures and empirically adequate clock/scale relations, it may describe local kinematics using spatial and temporal coordinates. The current cut construction has not derived those structures, their dimensionality, an Einstein metric, or proper time from causal order alone.
\end{proposition}

\begin{remark}[the open joint]{\normalfont\ \conj/\gap}
The missing bridge is from a declared influence/reachability structure and calibrated physical clocks to a tested Lorentzian geometry. The next section studies a model-dependent anisotropy obstruction; its result neither closes that bridge nor upgrades an entropic diagnostic to the source of order.
\end{remark}

\section{Requirements for emergent Lorentz invariance \quad{\normalsize\est\ / \conj}\label{sec:lorentz}}

A $3{+}1$ spacetime model would require Lorentzian effective kinematics beyond the causal-accessibility proposal of \S\ref{sec:spacetime}. This section lists the required components and states which are assumed, established in restricted models, or still missing. It does not derive Lorentz invariance.

\paragraph{The five components.}
Lorentz invariance bundles five distinct claims:
\begin{enumerate}[leftmargin=1.6em, itemsep=1pt, topsep=2pt]
  \item[\textbf{C1}] \textbf{An invariant finite speed} $c$ exists.
  \item[\textbf{C2}] \textbf{Relativity principle:} no frame is privileged.
  \item[\textbf{C3}] \textbf{Boost mixing:} space and time transform into one another under changes of velocity.
  \item[\textbf{C4}] \textbf{Invariant interval} with Lorentzian ($-,+,+,+$) signature.
  \item[\textbf{C5}] \textbf{Causal structure:} timelike / spacelike / null.
\end{enumerate}

\paragraph{What the framework already supplies.}
\begin{itemize}[leftmargin=1.4em, itemsep=2pt, topsep=2pt]
  \item \textbf{A relation to C2 is not established. \conj} ``No privileged cut'' (Claim~\ref{cl:noselect}) is an underdetermination claim; the relativity principle is a physical covariance claim about inertial frames and laws. They are not the same theorem, and relating them requires additional argument.
  \item \textbf{C5 requires a separate model. \asm/\conj} The influence relation $\leadsto_M$ asks whether admissible intervention variation changes an accessible downstream record distribution. Local interactions and bounds on propagation can constrain possible edges, but $\hat{H}_{AB}\neq 0$ alone proves neither a particular directed influence nor relativistic timelike/null structure. Relating this reachability to Lorentzian causal cones requires further physical assumptions and tests.
  \item \textbf{C1 has a candidate. \est} A finite maximum rate of correlation propagation exists in any system with local interactions: the \textbf{Lieb--Robinson bound}. This is a proven theorem, and it supplies an emergent, finite, \emph{pre-geometric} speed scale --- the fastest one region can come to be correlated with another. This is \emph{not} the same object as the observed relativistic conversion constant of \S\ref{sec:reach} (Route~4); it is a model-dependent effective speed. Whether the two coincide (i.e.\ whether the emergent correlation speed \emph{is} the observed $c$) is part of what the isotropy problem below must settle. Its numerical value is in any case a unit convention.
  \item \textbf{C3, C4 follow only with extra assumptions. \conj} Given an invariant finite speed (C1) and the relativity principle (C2), boost-mixing and the invariant interval follow \emph{only} with homogeneity, isotropy, linearity, and suitable group structure. Since isotropy is precisely the open obstruction below, this step cannot be leaned on as established; it is downgraded to conditional.
\end{itemize}

\paragraph{Anisotropy limitation. \gap}
The Lieb--Robinson speed is generally \textbf{anisotropic and frame-dependent} on a structured interaction pattern. A finite propagation bound therefore does not by itself produce Lorentz invariance. The effective propagation law would also need the required isotropy, boost covariance, and interval structure in the stated domain.

\paragraph{Ruler-cancellation proposal (Route C). \asm/\conj}
The proposal assumes that an observer's rulers and clocks are built from the interaction structure being measured. Cancellation requires a strong universality condition: rods, clocks, signals, and comparison procedures must inherit the same effective directional scaling. Endogenous instruments alone do not guarantee this condition; birefringence, dispersion, mode splitting, or orientation-dependent decay can expose anisotropy. Under the stated universality condition:

\begin{enumerate}[leftmargin=1.6em, itemsep=2pt, topsep=2pt]
  \item To measure a speed in direction $d$, an observer $A$ needs a rod and a clock, both physical systems built from $A$'s own couplings.
  \item \emph{If} rod, clock, and the measured process all share one effective dispersion, their common directional dependence enters every measurement in the same way.
  \item The shared directional factor then cancels in the dimensionless ratio that constitutes a speed measurement, so $A$ measures the same speed in every direction.
  \item Hence anisotropy is unregistrable --- \emph{provided} the co-dispersion is universal. If different sectors retain different anisotropies, the preferred structure becomes observable and the cancellation fails.
\end{enumerate}

\begin{proposition}[postulate about inaccessible anisotropy]{\normalfont\ \post}\label{prop:unreg}
The ontology may classify a difference as a physical fact only when some modeled access class can register it (\S\ref{sec:info}). Under that postulate, exact universal co-dispersion would make a preferred frame operationally absent. This is not a physics theorem and is not established here. Exactness would require every dimensionless Lorentz-violating observable to cancel for every relevant endogenous apparatus; the present model does not show that.
\end{proposition}

\paragraph{Possible relation to renormalization-group scaling. \conj}
A renormalization-group fixed point may give distinct low-energy observables a common effective dispersion. If the same scaling governs rods, clocks, and signals, the ruler-cancellation condition can hold in that regime.

\begin{proposition}[conditional relation between Routes B and C]{\normalfont\ \conj}\label{prop:routes}
If RG flow produces a common effective dispersion for the measured process and the observer's standards, then RG suppression of microscopic anisotropy (Route B) can also support ruler cancellation (Route C). The relation is conditional; away from that regime, residual anisotropy may remain observable.
\end{proposition}

\begin{remark}[a testable consequence]{\normalfont\ \conj}
RG suppression alone can leave residual Lorentz-violating corrections. Shared scaling in the observer's standards could suppress some dimensionless measurements further. The numerical experiment below tests this distinction for specific observables.
\end{remark}

\begin{remark}[status and limitations]{\normalfont\ \conj/\gap}
(i)~The framework does not show that relevant systems generically reach a fixed point or that static and dynamic dispersions coincide exactly. (ii)~The observer and apparatus definitions are model assumptions and cannot be chosen only to force cancellation. (iii)~An endogenous experiment requires the local sequencing and recordability already assumed by OC-A2--A3. Connecting that order to a thermodynamic arrow requires additional assumptions. Exact cancellation and its relation to macroscopic arrows remain conjectural.
\end{remark}

\paragraph{Numerical test. \est-in-model / \conj}
The ruler-cancellation claim was tested in a simulation of 2D anisotropic tight-binding free fermions ($J_x\neq J_y$) with exact correlation-matrix dynamics. The model-level control measures the Lieb--Robinson cone in bare lattice units. The endogenous observer uses impurity bound-state widths as rods and a cavity beat period as a clock, both built from the same couplings. The detection battery was first checked to detect anisotropy in model-level data and remain null on an isotropic substrate. The preregistered rules gave the following results.

\begin{itemize}[leftmargin=1.4em, itemsep=3pt, topsep=3pt]
  \item \textbf{Naive claim falsified.} An endogenous observer timing its \emph{fastest} signal (the Lieb--Robinson front) against its own rods measures a direction-dependent speed that plateaus at a non-cancelled value in the large-rod / near-fixed-point regime. ``No dimensionless internal experiment detects anisotropy'' is \textbf{false}. The battery detects the substrate axis.
  \item \textbf{Observed continuum scaling law.} Ballistic front speeds scale as $J_d$ (band-edge, UV); binding-length rods scale as $\sqrt{J_d}$. The measured anisotropy therefore approaches a square-root relation rather than cancellation:
  \[
    A_{\text{endo}}(\text{front}) \to \sqrt{A_{\text{god}}}
    \qquad(\text{measured } 1.4195 \text{ vs } \sqrt{2.0075}=1.4169).
  \]
  \item \textbf{Corrected claim supported, but IR-only.} Long-wavelength observables co-disperse with the rods and become isotropic at the fixed point ($A_{\text{endo}}(\text{wp})\to 1.0002$ as the bound-state size $\xi\to 39$; spreading central values move the same way but are tail-limited at the largest rods). This is the co-dispersion mechanism, numerically verified. But it holds for the \textbf{IR sector only}: the front (UV) observable never renormalizes.
\end{itemize}

\noindent The test supports cancellation only for the stated infrared observables. An endogenous observer with access to the tested ultraviolet front still detects substrate anisotropy. Whether an interacting fixed point changes the front scaling is outside this experiment and remains open.

\section{The observer, minimally defined \quad{\normalsize\asm}\label{sec:observer}}

The ruler-cancellation proposal (\S\ref{sec:lorentz}) requires a physical observer model. The definition below specifies only the subsystem, recording, and endogenous-standard assumptions used in that proposal. It does not model consciousness.

\begin{definition}[minimal observer]\label{def:observer}
A subsystem $A$ of a cut is an \textbf{observer} iff:
\begin{enumerate}[leftmargin=1.6em, itemsep=1pt, topsep=2pt]
  \item[\textbf{O1}] \textbf{It is one side of a declared cut.} The partition-relative quantities defined earlier can then be assigned. This condition alone is not sufficient.
  \item[\textbf{O2}] \textbf{It registers.} It has internal degrees of freedom whose states become correlated with $B$ so that relevant differences leave retained traces in $A$.
  \item[\textbf{O3}] \textbf{Its rulers are endogenous.} Its clocks and length-standards are built from the same $\hat{H}_{AB}$-type couplings that govern the propagation it measures. It has no structure-independent ruler.
\end{enumerate}
\end{definition}

\noindent Consciousness and a world model are not required. O2 supplies retained records; O3 supplies the endogenous standards used in \S\ref{sec:lorentz}.

\begin{remark}[the observer is downstream of the arrow]{\normalfont\ \gap}
O2 presupposes enough local ordering and trace stability for a record relation, already acknowledged in OC-A2--A3. It does not by itself require a universal thermodynamic arrow. Any stronger connection between record maintenance, a low-entropy boundary condition, and emergent Lorentz claims belongs to the specific realization and must be modelled rather than defined into ``observer.''
\end{remark}

\begin{remark}[scope of the observer definition]{\normalfont\ \asm}
The cancellation proposal concerns physical standards, not cognition. A richer observer model may be needed for other questions, but it is not required for this particular measurement model.
\end{remark}

\section{What the present realization does not select \quad{\normalsize\gap}\label{sec:select}}

Everything in \S2--\S4 presupposes a chosen cut: $\mathcal{H}_A\otimes\mathcal{H}_B$ in \S2--\S3, $K_B$ in \S4. The relevant question is what follows from the declared data $(\mathcal{H}_U,\lvert\Psi\rangle,\hat H_U)$ alone. Operational indistinguishability does not by itself prove that no physical cut-selection principle can exist; the following claim is restricted to this input class and remains \asm/\conj, not a general no-go theorem.


\begin{claim}[Selection from the present inputs]\label{cl:noselect}{\normalfont\ \asm/\conj}
No \textbf{basis-independent, unitary-equivariant, physically natural} selection rule has been supplied here that takes only the declared global Hilbert space, state, and Hamiltonian as input and outputs a preferred non-trivial factorization. Tensor factorizations are additional data relative to this formulation: for composite finite dimension, $\mathcal{H}_U$ admits many; for prime dimension no non-trivial factorization of this kind exists, and the infinite-dimensional case differs. A richer physical theory may add locality, an observable algebra, interaction structure, or other principles. Whether such structure genuinely selects a cut is a separate claim with separate assumptions, not ruled out by operational equivalence alone.
\end{claim}

\noindent\emph{Support (not full proof).} This is closely related to known results:
\begin{itemize}[leftmargin=1.4em, itemsep=1pt, topsep=2pt]
  \item The tensor-product structure is additional data, not fixed by $\mathcal{H}_U$ (factorization ambiguity / basis-dependence of entanglement).
  \item Candidate rules based on preferred observables, locality, interaction form, or an observer use structure beyond the present tuple. Such structure may be physically justified, but must be declared rather than counted as an output of the tuple alone.
\end{itemize}

\begin{corollary}[formal counterpart of the regress]\label{cor:regress}{\normalfont\ \asm}
Within the present input set, an actual selection refers to further structure (preferred observables, an observer, or a locality assignment) not furnished by the tuple itself. Calling that structure ``exogenous'' identifies a missing input in this realization; it does not prove an infinite regress for every possible physical theory.
\end{corollary}

\begin{remark}[requirement on added selection structure]
A realization that selects a cut through added locality or observable structure must declare and test that addition. This may be legitimate physical explanation rather than mere ``smuggling''; the point of Claim~\ref{cl:noselect} is narrower: the present tuple does not supply such a rule, so the chosen factorization cannot be advertised as an output of this model.
\end{remark}

\section{Two unresolved inputs \quad{\normalsize\texorpdfstring{\gap\ $\times 2$}{[GAP] x 2}}\label{sec:turtles}}

The current construction names two unresolved inputs, historically called ``the two turtles.'' This is bookkeeping for this realization, not a proof that every possible theory has exactly two foundational assumptions.

\paragraph{Arrow conditions. \gap}
$S_A(t)$ is not generically monotone under unitary evolution (the direction remark in \S\ref{sec:time}). The present interpretation proposes a low-entropy boundary condition on $\lvert\Psi\rangle$ (the past hypothesis), plus a relevant coarse-graining and environmental account. A low global entropy condition alone does not guarantee monotonic reduced entropy for every cut. This input is
\begin{itemize}[leftmargin=1.4em, itemsep=1pt, topsep=2pt]
  \item \textbf{not derivable} from $(\mathcal{H}_U, \hat{H}_U)$ alone \est{} (open problem);
  \item a constraint on the \emph{state}, entering as an initial condition;
  \item part of the additional structure needed for the directionality claims made in this realization.
\end{itemize}
Formally: the arrow is an \textbf{external boundary condition} --- $\langle\Psi_0\vert \hat{S}_{\text{coarse}}\vert\Psi_0\rangle$ low for a chosen coarse-graining --- not a theorem.

\paragraph{Cut selection. \gap}
By Claim~\ref{cl:noselect}, the present input tuple does not select the cut. This is
\begin{itemize}[leftmargin=1.4em, itemsep=1pt, topsep=2pt]
  \item \textbf{not a state constraint} but a \emph{selection}, made from a vantage;
  \item not supplied by the arrow condition --- fixing the initial state's entropy does not fix the factorization.
\end{itemize}

\begin{proposition}[distinct roles in the current bookkeeping]\label{prop:indep}{\normalfont\ \asm}
As represented here, the arrow input includes a boundary/preparation condition, while the cut input supplies a factorization. Neither datum, as presently defined, contains a rule returning the other. They are therefore tracked as separate unresolved inputs in this construction. This bookkeeping claim does not establish logical independence across every richer theory that may relate state, dynamics, locality, and factorization.
\end{proposition}

\noindent This rejects collapsing the two questions without an explicit bridge; it does not prove that all possible theories must bottom out in exactly two assumptions.

\section{A conditional route to intersubjectivity \quad{\normalsize\conj\ / \gap}\label{sec:inter}}

Because the present realization does not select cuts, it does not guarantee that observers with distinct cuts $(A\vert B)$ and $(A'\vert B')$ register compatible records. This section gives a conditional mechanism based on redundant records. It does not solve intersubjectivity or reduce its assumptions to the two unresolved inputs above.

\paragraph{The reduction.}
\begin{enumerate}[leftmargin=1.6em, itemsep=2pt, topsep=3pt]
  \item \textbf{Redundancy can support shareability.} \est{} (quantum Darwinism, Zurek). In the relevant decoherence models, many observers can retrieve compatible information from distinct environment fragments when records are redundantly proliferated. This is a sufficient mechanism in that setting, not a universal iff-definition of intersubjectivity.
  \item \textbf{Redundant proliferation requires a realization.} \conj{} A particular decoherence model may connect environmental preparation, record proliferation, and an arrow. The connection requires declared boundary conditions, correlations, coarse-graining, and readout resources; it is not an identity between entropy, records, and time.
\end{enumerate}

\begin{proposition}[conditional shareability mechanism]\label{prop:inter}{\normalfont\ \asm/\conj}
\emph{Given} pointer-aligned access, an appropriate environment, and redundant, stable records, distinct observers may retrieve mutually consistent facts from different fragments. Quantum Darwinism motivates this route, but the required observer population, access map, and stability conditions are additional assumptions. Intersubjectivity remains open beyond that conditional domain.
\end{proposition}

\paragraph{Residual dependencies. \gap}
Proposition~\ref{prop:inter} assumes pointer-aligned operational access. Quantum Darwinism can identify observables whose records proliferate in a specified interaction model; it does not by itself select an observer's factorization, apparatus, or access class. The agreement mechanism is therefore conditional on additional cut/access alignment and cannot yet be used to reduce intersubjectivity to the arrow.

\paragraph{Observer concentration is an open model.}
One may hypothesize that an observer ensemble generated by the same dynamics concentrates near record-compatible cuts, but neither the ensemble nor a measure over its cuts is defined here. The claim that off-pointer observers are rare is therefore not established.

\begin{remark}[a possible Boltzmann-type ansatz]{\normalfont\ \conj}
A weighting such as $\exp\!\big(-\mathcal{C}_{\text{dec}}/\lambda\big)$ is only an ansatz until the factorization/access space, measure, decoherence cost $\mathcal{C}_{\text{dec}}$, scale $\lambda$, and observer-population ensemble are specified. Its local or global shape cannot be inferred from the slogan alone. This revision retains the proposal as future work and makes no concentration theorem.
\end{remark}

\begin{corollary}[conditional distribution proposal\label{cor:dist}]{\normalfont\ \asm}
If a well-defined observer ensemble, measure on cuts, and decoherence cost are supplied, the boundary/preparation conditions and interaction model may jointly shape a distribution over record-bearing observers. Those ingredients are not defined here, so this does not establish that intersubjectivity introduces no additional foundational assumptions.
\end{corollary}

\begin{remark}[remaining open questions]{\normalfont\ \conj/\gap}
Two links above are physically motivated but not proved here: (i)~the conditions under which redundancy tracks a macroscopic arrow, and (ii)~whether any observer concentration has the proposed Boltzmann form. Making the second rigorous requires a measure on the relevant factorization/access space, a defined observer ensemble, and a derivation of the weighting. Intersubjectivity therefore remains a named open problem; the present sketch does not reduce it to the other two gaps.
\end{remark}

\section{What this spec is and is not}\label{sec:isnot}

\noindent\textbf{Is} \asm:
\begin{itemize}[leftmargin=1.4em, itemsep=1pt, topsep=2pt]
  \item a proposed operational core plus a quantum realization with a conditional causal-accessibility relation, clock/arrow distinctions, reduced entropy (\S2), Noether (\S3), Shannon quantities and Fields' GHP (\S3--\S4), and an observer-registration condition (\S4);
  \item an explicit scope ledger: the present input tuple does not supply a cut selector; arrow and cut assumptions are tracked separately here; and the suggested route to intersubjectivity remains conditional.
\end{itemize}

\noindent\textbf{Is not:}
\begin{itemize}[leftmargin=1.4em, itemsep=1pt, topsep=2pt]
  \item a derivation of Lorentz invariance, a metric, proper time, or an objective global time from causal order alone;
  \item a derivation of the arrow (the required boundary and preparation conditions are additional inputs, \S\ref{sec:turtles});
  \item a selection of the cut (Claim~\ref{cl:noselect} says this realization does not supply one, without ruling out every richer physical principle);
  \item a \emph{closed} solution to intersubjectivity (\S\ref{sec:inter} offers only a conditional mechanism with named assumptions);
  \item new mathematics: every \est/\fields{} piece is prior work; the contribution is arrangement plus the explicit gap-statements.
\end{itemize}

\noindent\textbf{The proposed organizing claims, if developed:} (i)~the current realization tracks arrow conditions and cut selection as distinct inputs unless a richer theory supplies a bridge; (ii)~redundant-record dynamics may explain a restricted form of intersubjective agreement once an access map, environment, observer ensemble, and measure are specified. These are research proposals with explicit failure conditions, not established facts about every possible theory.

\section{Immediate next steps (if pursued)}\label{sec:next}

\begin{enumerate}[leftmargin=1.6em, itemsep=2pt, topsep=3pt]
  \item \textbf{Scope the no-selection claim.} Compare Claim~\ref{cl:noselect} with the literature and state the input class covered by each result. Do not extend the claim beyond those assumptions.
  \item \textbf{Specify bridges} among state, dynamics, locality, and factorization. Do this before claiming logical independence or that selection is impossible.
  \item \textbf{Make \S\ref{sec:inter} testable.} Define the access maps, observer ensemble, measure on the relevant cut space, decoherence cost, and record-stability criteria. Only then test whether the proposed concentration and shareability claims hold.
\end{enumerate}

\paragraph{Future causal-reconstruction pilot (not implemented).}
A later brickwork quantum circuit or cellular automaton could supply a strict finite model-level causal cone while hiding simulator steps and exact lattice coordinates from its observer. Tagged interventions and detector records would support a reconstruction $\widehat C_O$ of $C_M=(E,\prec_M)$. Access could expand in stages: restricted detectors; added detectors and references; an endogenous clock; then an asymmetric environment with persistent records. Only a separately specified experiment should define edge-detection, false-positive, missed-edge, reachability, cone-boundary, or model-uncertainty metrics. No such pilot or metric is implemented in this revision.

\end{document}
